\section{The completed mock forms and their Mellin equations}
\label{sec:mock-mellin}
The completed mock forms have a modular functional equation, but their
holomorphic parts need not have an ordinarily convergent Dirichlet series.
The appropriate general framework uses test functions, as in Diamantis,
Lee, Raji and Rolen \cite[Definition~4.2, Lemma~4.4 and Theorem~4.5]{DLRR2023}.
We give both that equation and a specified regularized Mellin version.

Let $\widehat H$ be an umbral completion of weight $1/2$, with the odd
component convention $H_{-r}=-H_r$.  Write $G=\xi_{1/2}\widehat H$, where
\[
 \xi_kF=2iy^k\overline{\frac{\partial F}{\partial\bar\tau}}.
\]
This defines the normalization of $G$; it is a fixed scalar multiple of the
unary shadow $S^X$.  With this convention the nonholomorphic part is
\begin{equation}\label{eq:mock-nonholomorphic-part}
 H^-(\tau)=-\sum_{\nu>0}\overline{b(\nu)}(4\pi\nu)^{-1/2}
       \Gamma(1/2,4\pi\nu y)q^{-\nu},
 \qquad G(\tau)=\sum_{\nu>0}b(\nu)q^\nu.
\end{equation}
The values of $\nu$ may be rational.  The normalization is useful because
applying $\xi_{1/2}$ to each summand gives exactly $b(\nu)q^\nu$.

For the untwined form, put $N=2m$ and define, on the odd subspace,
\begin{equation}\label{eq:mock-weil-involution}
 W_{r,r'}=\frac{i}{\sqrt N}e^{2\pi irr'/N}.
\end{equation}
The Jacobi transformation law and theta decomposition give
\begin{equation}\label{eq:mock-axis-transformation}
 \widehat H(i/y)=y^{1/2}W\widehat H(iy),\qquad W^2=I.
\end{equation}
Indeed the Fourier matrix squared is $N$ times $r\mapsto-r$, which is
$-I$ on odd vectors; the factor $i^2=-1$ gives $W^2=I$.
At $m=6$, equation~\eqref{eq:mock-weil-involution} is the explicit
$12\times12$ formula $W_{r,r'}=i e^{2\pi irr'/12}/\sqrt{12}$,
restricted to $H_0=H_6=0$ and $H_{12-r}=-H_r$.
The same argument applies at every lambency.  For a twining whose modular
group has several relevant cusps, one uses the vector of cusp expansions
and the corresponding matrix in place of $W$.

\begin{proposition}[Test-function equation]\label{prop:mock-test-function}
For a smooth compactly supported function $\phi:(0,\infty)\to\mathbb C$, put
\[
 \mathcal L_{\widehat H}(\phi)=
       \int_0^\infty\widehat H(iy)\phi(y)\,dy,
 \qquad \phi^\vee(y)=y^{-3/2}\phi(1/y).
\]
Then
\begin{equation}\label{eq:mock-test-functional}
 \mathcal L_{\widehat H}(\phi)
       =W\mathcal L_{\widehat H}(\phi^\vee).
\end{equation}
\end{proposition}
\begin{proof}
All integrals converge because the support is a compact subset of
$(0,\infty)$.  Substitute $y=1/t$ and use
\eqref{eq:mock-axis-transformation}.  The Jacobian $t^{-2}$ combines with
$t^{1/2}$ to give $t^{-3/2}$, which is precisely the displayed dual test
function.
\end{proof}

A plain power $\phi(y)=y^{s-3/2}$ is not an admissible test function when
$\widehat H$ has a polar Fourier term: that term grows exponentially as
$y\to\infty$.  To make the Mellin notation unambiguous, fix the following
regularization.  Write the finitely many growing and constant terms at
infinity as
$P(y)=\sum_{\nu<0}a(\nu)e^{-2\pi\nu y}+a(0)$ and put
\begin{align}
 F_H(u)={}&\int_1^\infty\bigl(\widehat H(iy)-P(y)\bigr)y^{u-1}\,dy
 -\frac{a(0)}u\notag\\
 &+\sum_{\nu<0}a(\nu)(2\pi\nu+i0)^{-u}
                 \Gamma(u,2\pi\nu+i0).\label{eq:mock-tail-regularization}
\end{align}
Here $+i0$ specifies continuation from the upper half-plane and thereby
fixes the branch on the negative real axis.  The integral is entire in
$u$ because its integrand decays exponentially at infinity; the other
terms are meromorphic.  Define the shifted, symmetrically regularized
Mellin transform by
\begin{equation}\label{eq:mock-regularized-Mellin}
 \mathcal M_{\widehat H}(s)=F_H(s-1/2)+WF_H(1-s).
\end{equation}
This is the regularized form of
$\int_0^\infty\widehat H(iy)y^{s-3/2}\,dy$ obtained by splitting at $1$
and using the modular transformation on the lower interval.
Its functional equation is now precise:
\begin{equation}\label{eq:mock-Mellin-FE}
 \mathcal M_{\widehat H}(s)
       =W\mathcal M_{\widehat H}(3/2-s).
\end{equation}
It follows immediately from $W^2=I$ and
\eqref{eq:mock-regularized-Mellin}.

\begin{proposition}[The shadow contribution]\label{prop:mock-shadow-Mellin}
Let $\overline G(\tau)=\sum_{\nu>0}\overline{b(\nu)}q^\nu$ and set
\begin{equation}\label{eq:mock-kernel}
 C(s)=(4\pi)^{-s}\frac{\Gamma(s)}{s-1/2}
 {}_2F_1\left(s-\tfrac12,s;s+\tfrac12;\tfrac12\right).
\end{equation}
For $\Re s>1$, the ordinary Mellin transform of $H^-$ converges absolutely
and is
\begin{equation}\label{eq:mock-shadow-Mellin}
 \int_0^\infty H^-(iy)y^{s-3/2}\,dy=-C(s)L(\overline G,s).
\end{equation}
Let $\mathcal M_{H^+}(s)$ denote the holomorphic-part Mellin transform with
the regularization fixed by
\eqref{eq:mock-tail-regularization}--\eqref{eq:mock-regularized-Mellin} and
linearity.  Then, by meromorphic continuation,
\begin{align}
 \mathcal M_{H^+}(s)-C(s)L(\overline G,s)
 =W\bigl\{\mathcal M_{H^+}(3/2-s)
       -C(3/2-s)L(\overline G,3/2-s)\bigr\}.
 \label{eq:mock-inhomogeneous-Mellin}
\end{align}
Every shadow $L$-function here is explicitly a sum of the
$L(2s-1,\psi)$ in Theorem~\ref{thm:all-shadow-characters}.
\end{proposition}
\begin{proof}
For a single term of \eqref{eq:mock-nonholomorphic-part}, set $t=4\pi\nu y$.
The resulting scalar integral, without the coefficient
$-\overline{b(\nu)}(4\pi\nu)^{-s}$, is
\[
 K(s)=\int_0^\infty t^{s-3/2}e^{t/2}\Gamma(1/2,t)\,dt.
\]
For $\Re s>1/2$, insert
$\Gamma(1/2,t)=\int_t^\infty e^{-v}v^{-1/2}\,dv$, expand $e^{t/2}$, and
integrate first over $0<t<v$.  Absolute convergence gives
\begin{align*}
 K(s)&=\sum_{j\geq0}\frac{\Gamma(s+j)}{2^j j!(s+j-1/2)}\\
 &=\frac{\Gamma(s)}{s-1/2}
 {}_2F_1\left(s-\tfrac12,s;s+\tfrac12;\tfrac12\right).
\end{align*}
For a unary shadow, $\sum_{\nu>0}|b(\nu)|\nu^{-\sigma}$ converges for
$\sigma>1$, so summation and integration are justified in the stated
half-plane.  This proves \eqref{eq:mock-shadow-Mellin}.
Subtract this convergent term from the fixed regularized transform of
$\widehat H$ and continue meromorphically.  Substitution into
\eqref{eq:mock-Mellin-FE} proves
\eqref{eq:mock-inhomogeneous-Mellin}.
\end{proof}

The Dirichlet Riemann hypothesis therefore controls specified constituents
of the inhomogeneous shadow term.  It gives no zero-location theorem for
$\mathcal M_{H^+}$ or for $\mathcal M_{\widehat H}$ by this identity.
Moreover, the naive expression $\sum a(\nu)\nu^{-s}$ for a holomorphic
mock part may fail to converge in every half-plane because polar modular
forms can have coefficients of exponential order in $\sqrt\nu$.
A question about its zeros must first specify a regularization or a
test-function family.  The choice in
\eqref{eq:mock-tail-regularization} is explicit; it is not an Euler product,
and equation~\eqref{eq:mock-inhomogeneous-Mellin} provides no argument that
its zeros lie on a vertical line.
